\chapter*{ABSTRACT}
\phantomsection
\addcontentsline{toc}{chapter}{ABSTRACT}

\textbf{Background.} Human papillomavirus (HPV) vaccination is a highly effective primary prevention strategy for cervical cancer, yet vaccine awareness and uptake remain unequally distributed in many low- and middle-income settings. National evidence on socioeconomic inequality in HPV vaccine awareness, self-reported any-dose uptake, and the awareness--vaccination gap among young Vietnamese women is limited.

\textbf{Objectives.} This thesis aimed to estimate the prevalence of HPV vaccine awareness, self-reported any-dose HPV vaccine uptake, and the awareness--vaccination gap among Vietnamese women aged 15--29 years, assess associated sociodemographic factors, quantify wealth-related inequality using concentration indices, and decompose the observed contributors to inequality.

\textbf{Methods.} A cross-sectional secondary analysis was conducted using the Viet Nam Multiple Indicator Cluster Survey 2020--2021. The analytic sample included 4,102 women aged 15--29 years with valid survey design information. All descriptive and bivariate analyses accounted for sampling weights, stratification, and clustering. Bivariate p-values were obtained from design-based Pearson chi-squared tests with survey correction. Vaccination was measured as self-reported receipt of any HPV vaccine dose, not schedule completion. Adjusted prevalence ratios (aPRs) were estimated using survey-weighted Poisson regression. Wealth-related inequality was measured using the standard concentration index and the Erreygers-corrected concentration index; decomposition used probit average marginal effects.

\textbf{Results.} Overall, 62.4\% of women had heard of HPV vaccination (95\% CI: 59.8--64.9), while 7.5\% reported any-dose HPV vaccine uptake (95\% CI: 6.2--8.7). Among women who were aware of HPV vaccination, only 12.0\% were vaccinated (95\% CI: 10.1--13.9), leaving an awareness--vaccination gap of 88.0\% (95\% CI: 86.1--89.9). Both awareness and vaccination were concentrated among wealthier women. The Erreygers concentration index was 0.371 (95\% CI: 0.328--0.415) for awareness and 0.116 (95\% CI: 0.089--0.144) for vaccination. The gap was concentrated among poorer aware women (Erreygers CI: -0.123; 95\% CI: -0.164 to -0.082). In adjusted models, the poorest quintile had substantially lower awareness (aPR 0.576; 95\% CI: 0.480--0.691) and vaccination (aPR 0.118; 95\% CI: 0.039--0.358) than the richest quintile, and a higher prevalence of remaining unvaccinated despite awareness (aPR 1.137; 95\% CI: 1.055--1.225).

\textbf{Conclusion.} HPV vaccine awareness is moderately common among young Vietnamese women, but self-reported any-dose uptake remains low. Socioeconomic position, education, residence, media exposure, ethnicity, and gender-equitable attitudes were associated with different stages of the awareness-to-vaccination pathway. These estimates provide a pre-program equity baseline for monitoring any-dose HPV vaccine uptake as public financing is expanded. Equity-oriented HPV vaccination policy should be designed to reduce inequity by combining public financing with targeted outreach for poorer, rural, less educated, and ethnic minority communities.

\textbf{Keywords:} HPV vaccination; cervical cancer prevention; socioeconomic inequality; concentration index; Erreygers decomposition; MICS; Viet Nam.
\clearpage
